\section{Introdução}
\label{sec:introducao}

Em carregadores de veículos elétricos o estágio de entrada precisa absorver uma corrente
senoidal com FP elevado e baixo THD, dentro das normas estabelecidas pela IEEE 519 \cite{ieee519}
e fornecer um barramento CC regulado ao conversor CC-CC que alimenta a bateria. Os conversores PFC monofásicos, 
como o \textit{boost}, o \textit{bridgeless} e o \textit{totem-pole}, são amplamente utilizados em carregadores de baixa potência,
porém em estações de recarga rápidas ou carregadores embarcados que operam na faixa de \SI{20}{kW}, faz-se necessária a utilização
de retificadores trifásicos multiníveis, como o retificador de Vienna, para reduzir os esforços nas chaves semicondutoras e o tamanho de filtros utilizados \cite{yilmaz2013review,tu2019extreme}.

O retificador de Vienna se destaca nesse cenário por utilizar apenas três chaves bidirecionais, que bloqueiam metade da tensão do barramento, 
dessa forma diminuindo o estresse nas chaves e permitindo a utilização de semicondutores de menor tensão nominal \cite{safayatullah2022comprehensive}, sendo amplamente utilizado em módulos comerciais de recarga rápida.

O retificador de Vienna a quatro fios (ponto de neutro aterrado) com fator de potência variável é uma topologia em que o neutro da rede é ligado ao ponto médio do barramento \cite{alahuhtala2006fourwire}, essa conexão permite
que o conversor opere com cargas e redes desequilibradas, tendo em vista que ele fixa o potencial do ponto médio (o que reduz o ruído de modo comum) e viabiliza sua utilização em redes monofásicas e trifásicas \cite{schafmeister2021transformerless}. 
Além disso, retificadores unidirecionais, além de poderem operar com PFC, podem oferecer serviços à rede como compensação de reativo e FP ajustável. 
Nesse sentido, este trabalho visa o desenvolvimento de uma estrutura de controle de um retificador de Vienna trifásico a quatro fios no referencial estacionário, com PR por fase e FP variável desacoplado da malha de tensão.
